
\documentclass{article}
\usepackage{colortbl}
\usepackage{makecell}
\usepackage{multirow}
\usepackage{supertabular}

\begin{document}

\newcounter{utterance}

\centering \large Interaction Transcript for game `codenames', experiment `ambiguous', episode 1 with reflexion{-}Qwen3.5{-}27B.
\vspace{24pt}

{ \footnotesize  \setcounter{utterance}{1}
\setlength{\tabcolsep}{0pt}
\begin{supertabular}{c@{$\;$}|p{.15\linewidth}@{}p{.15\linewidth}p{.15\linewidth}p{.15\linewidth}p{.15\linewidth}p{.15\linewidth}}
    \# & $\;$A & \multicolumn{4}{c}{Game Master} & $\;\:$B\\
    \hline

    \theutterance \stepcounter{utterance}  
    & & & \multicolumn{2}{p{0.3\linewidth}}{
        \cellcolor[rgb]{0.9,0.9,0.9}{
            \makecell[{{p{\linewidth}}}]{
                \texttt{\tiny{[GM$|$GM]}}
                {'hidden': {'team': ['dab', 'flock'], 'innocent': ['steer', 'darn', 'stir', 'counter'], 'opponent': ['bust'], 'assassin': []}, 'revealed': {'team': {'team': [], 'innocent': [], 'opponent': [], 'assassin': []}, 'opponent': {'team': [], 'innocent': [], 'opponent': [], 'assassin': []}}}
            }
        }
    }
    & & \\ \\

    \theutterance \stepcounter{utterance}  
    & & \multicolumn{4}{p{0.6\linewidth}}{
        \cellcolor[rgb]{0.9,0.9,0.9}{
            \makecell[{{p{\linewidth}}}]{
                \texttt{\tiny{[P1$\langle$GM]}}
                \texttt{You have attempted this same task before and did not succeed. These are the plans you wrote after those attempts.} \\
\texttt{Trial \#1: New plan: My previous attempt failed because I listed words in the "TARGETS" field that were not in my team's word list (e.g., "racket", "serve") and I generated an excessively long, repetitive list that violated the format constraints. I also failed to update the list based on the remaining words provided by the Game Master. Next time, I will strictly select target words only from the provided "remaining team words" list. I will ensure the clue is a single word that semantically links to a subset of these specific words without triggering opponent, distractor, or kill words. Finally, I will format the response exactly as "CLUE: <word>" followed by "TARGETS: <comma{-}separated list of valid team words only>", keeping the list concise and free of repetition.} \\
\\ 
\texttt{Trial \#2: New plan: My previous attempt failed because I listed words in the "TARGETS" field that were not in my team's current word list (e.g., "racket", "serve" instead of valid words like "match" or "jet"). I must strictly select target words only from the provided "remaining team words" list: sheet, match, dab, jet, pump, flock, rake, lime. I will ensure the clue is a single word semantically linked to a subset of these specific valid words, avoiding any opponent, distractor, or kill words. Finally, I will format the response exactly as "CLUE: <word>" followed by "TARGETS: <comma{-}separated list of valid team words only>", ensuring no extra text or invalid words are included.} \\
\\ 
\texttt{Trial \#3: New plan: My previous attempt failed because I chose clues that were too broad or ambiguous, causing my partner to guess distractor words (like "stalk" for "flock" and "pipe" for "rake") instead of the intended targets. Specifically, "GARDEN" linked to "rake" but also strongly suggested "pipe" (distractor), and "CLOTH" likely led to a dead end or incorrect guess for "dab". Next time, I will prioritize clues that have a unique, strong semantic link to my specific remaining words while having no obvious connection to the current distractor or opponent lists. For example, instead of "GARDEN", I should have used "TOOL" for "rake" (avoiding "pipe") or "BIRD" for "flock" (avoiding "stalk"). I will also avoid listing "dab" with "sheet" under "CLOTH" if "dab" is better associated with "PAINT" or "MAKEUP" to ensure distinctness. I will carefully cross{-}reference my chosen clue against the *current* list of distractors before submitting to prevent turn{-}ending mistakes.} \\
\\ 
\\ 
\texttt{Provide one single{-}word clue that relates to some of your team's words from the list below. You can choose to how many and to which words your clue relates to.} \\
\texttt{The goal is to have your partner guess all of your team's words based on the clues you give them, before the other team has uncovered all of theirs. If your partner guesses an 'opponent' or a 'distractor' word, their guessing turn ends and the opposing team gets to make a turn. If your partner guesses a 'kill' word, the game ends and your team loses immediately.} \\
\texttt{The more words your clue relates to, the quicker you can win the game, but the harder the associations may be to guess for your partner, so choose your clue word wisely. The clue word has to be semantically related to the target words, it cannot be one of the words in the lists or contain parts of them.} \\
\\ 
\texttt{Always give your single{-}word clue and your comma{-}separated list of related target words in the following format and make your answers as short as possible, never include any other text than is required in this form:} \\
\\ 
\texttt{CLUE: <WORD>} \\
\texttt{TARGETS: <TARGETS>} \\
\\ 
\texttt{Your team words are: dab, flock.} \\
\texttt{Your opponent's team words are: bust.} \\
\texttt{Distractor words are: steer, darn, stir, counter.} \\
\texttt{The kill words are: .} \\
            }
        }
    }
    & \\ \\

    \theutterance \stepcounter{utterance}  
    & \multicolumn{4}{p{0.6\linewidth}}{
        \cellcolor[rgb]{0.9,0.9,0.9}{
            \makecell[{{p{\linewidth}}}]{
                \texttt{\tiny{[P1$\rangle$GM]}}
                \texttt{CLUE: BIRD} \\
\texttt{TARGETS: flock} \\
            }
        }
    }
    & & \\ \\

    \theutterance \stepcounter{utterance}  
    & & & \multicolumn{2}{p{0.3\linewidth}}{
        \cellcolor[rgb]{0.9,0.9,0.9}{
            \makecell[{{p{\linewidth}}}]{
                \texttt{\tiny{[GM$|$GM]}}
                \texttt{bird} \\
            }
        }
    }
    & & \\ \\

    \theutterance \stepcounter{utterance}  
    & & & \multicolumn{2}{p{0.3\linewidth}}{
        \cellcolor[rgb]{0.9,0.9,0.9}{
            \makecell[{{p{\linewidth}}}]{
                \texttt{\tiny{[GM$|$GM]}}
                ['flock']
            }
        }
    }
    & & \\ \\

    \theutterance \stepcounter{utterance}  
    & & & \multicolumn{2}{p{0.3\linewidth}}{
        \cellcolor[rgb]{0.9,0.9,0.9}{
            \makecell[{{p{\linewidth}}}]{
                \texttt{\tiny{[GM$|$GM]}}
                {'word': 'flock', 'assignment': 'team'}
            }
        }
    }
    & & \\ \\

    \theutterance \stepcounter{utterance}  
    & & \multicolumn{4}{p{0.6\linewidth}}{
        \cellcolor[rgb]{0.9,0.9,0.9}{
            \makecell[{{p{\linewidth}}}]{
                \texttt{\tiny{[GM$\rangle$P2]}}
                \texttt{You have attempted this same task before and did not succeed. These are the plans you wrote after those attempts.} \\
\texttt{Trial \#1: New plan: The previous attempt was aborted because the response included the word "GUESS" followed by a colon and a space, but the instruction explicitly required the list to *start* with 'GUESS: ' and contain *no other text*. While the format seemed correct at first glance, the rejection suggests a strict parsing issue, possibly due to the specific spacing or the interpretation of "start your list... with 'GUESS: '". To ensure compliance, I will strictly output the string "GUESS: " immediately followed by the comma{-}separated words with no leading or trailing whitespace, ensuring the entire response is exactly "GUESS: peck, cock, hatch, flock" with no extra characters or formatting deviations. I will double{-}check that the words selected (peck, cock, hatch, flock) are indeed from the provided list and relate to 'bird'.} \\
\\ 
\texttt{Trial \#2: New plan: The previous attempt was aborted because I provided four words ("cock, hatch, flock" plus an implied fourth or a miscount in the logic, though the text shows three, the error likely stems from the strict "up to 3" constraint or a hidden formatting nuance). The instruction explicitly limits the list to "up to 3 words". My previous response listed three words, which should be valid, but the "ABORT" status suggests a strict violation. Re{-}reading the prompt, the list provided is: pump, match, pipe, steer, squash, counter, stalk, stir, darn, jam, sheet, lime, cock, rake, scarf, vault, scale, bust, hatch, wax, jet, flock, dab. The words "cock", "hatch", and "flock" are all in the list. The error might be that the system expects exactly 3 or fewer, but perhaps the parsing failed due to the specific words chosen or a hidden rule. However, looking at the "ABORT" reason in the example, it often relates to format. Wait, the previous attempt output "GUESS: cock, hatch, flock". This is 3 words. Why was it aborted? Perhaps the system considers "cock" inappropriate or the list parsing is strict about the order? No, the most likely issue in these specific "ABORT" scenarios with "up to 3" is that the model might have hallucinated a word not in the list or the system has a strict character limit. But looking closer at the prompt's "ABORT" definition: "violated the game rules... did a thing you were specifically asked not to do". The rule is "up to 3 words". I provided 3. Maybe the issue is the space after the colon? The prompt says "start your list of guess(es) with 'GUESS: '". My output was "GUESS: cock, hatch, flock". This matches. Is it possible the system rejected "cock" due to content filters? Or perhaps the "ABORT" was a false positive in the simulation, but I must assume a rule violation. Let's re{-}evaluate the "up to 3" constraint. If I provide 3, it is valid. If I provide 4, it is invalid. I provided 3. Maybe the issue is that I need to ensure the words are *strictly* from the list and I might have misread the list? "cock" is there. "hatch" is there. "flock" is there. } \\
\\ 
\texttt{Let's reconsider the "ABORT" reason in the context of the example. In Example 2, the abort was due to a 3{-}letter word instead of 5. Here, maybe the abort is because the response contained *only* the guess, but the system expected something else? No, "do not include any other text".} \\
\\ 
\texttt{Trial \#3: New plan: The previous attempt failed because I included "innocent" words (words not strongly associated with the target) in my guesses, such as "stalk" for 'green' and "pipe" for 'garden'. The game likely requires high{-}precision associations to reach the target state, and including weak links caused the loss. In the future, I will strictly select only the single most obvious and direct association from the provided list for each target word. If a second word is not immediately and strongly related, I will provide only one word to avoid diluting the answer with incorrect or weak associations. For example, for 'green', I would have chosen only "lime" instead of adding "stalk", and for 'garden', only "rake" instead of adding "pipe".} \\
\\ 
\\ 
\texttt{Provide a comma{-}separated list of up to 1 words from the following list that best relate or are most closely associated with the word 'bird'. Always start your list of guess(es) with 'GUESS: ' and do not include any other text in your answer.} \\
\\ 
\texttt{steer, counter, stir, darn, bust, flock, dab} \\
            }
        }
    }
    & \\ \\

    \theutterance \stepcounter{utterance}  
    & & & \multicolumn{4}{p{0.6\linewidth}}{
        \cellcolor[rgb]{0.9,0.9,0.9}{
            \makecell[{{p{\linewidth}}}]{
                \texttt{\tiny{[GM$\langle$P2]}}
                \texttt{GUESS: flock} \\
            }
        }
    }
     \\ \\

    \theutterance \stepcounter{utterance}  
    & & & \multicolumn{2}{p{0.3\linewidth}}{
        \cellcolor[rgb]{0.9,0.9,0.9}{
            \makecell[{{p{\linewidth}}}]{
                \texttt{\tiny{[GM$|$GM]}}
                ['flock']
            }
        }
    }
    & & \\ \\

    \theutterance \stepcounter{utterance}  
    & & & \multicolumn{2}{p{0.3\linewidth}}{
        \cellcolor[rgb]{0.9,0.9,0.9}{
            \makecell[{{p{\linewidth}}}]{
                \texttt{\tiny{[GM$|$GM]}}
                {'word': 'flock', 'assignment': 'team'}
            }
        }
    }
    & & \\ \\

    \theutterance \stepcounter{utterance}  
    & & & \multicolumn{2}{p{0.3\linewidth}}{
        \cellcolor[rgb]{0.9,0.9,0.9}{
            \makecell[{{p{\linewidth}}}]{
                \texttt{\tiny{[GM$|$GM]}}
                {'word': 'flock', 'assignment': 'team'}
            }
        }
    }
    & & \\ \\

    \theutterance \stepcounter{utterance}  
    & & & \multicolumn{2}{p{0.3\linewidth}}{
        \cellcolor[rgb]{0.9,0.9,0.9}{
            \makecell[{{p{\linewidth}}}]{
                \texttt{\tiny{[GM$|$GM]}}
                {'word': 'bust', 'assignment': 'opponent'}
            }
        }
    }
    & & \\ \\

    \theutterance \stepcounter{utterance}  
    & & & \multicolumn{2}{p{0.3\linewidth}}{
        \cellcolor[rgb]{0.9,0.9,0.9}{
            \makecell[{{p{\linewidth}}}]{
                \texttt{\tiny{[GM$|$GM]}}
                {'hidden': {'team': ['dab'], 'innocent': ['steer', 'darn', 'stir', 'counter'], 'opponent': [], 'assassin': []}, 'revealed': {'team': {'team': ['flock'], 'innocent': [], 'opponent': [], 'assassin': []}, 'opponent': {'team': [], 'innocent': [], 'opponent': ['bust'], 'assassin': []}}}
            }
        }
    }
    & & \\ \\

    \theutterance \stepcounter{utterance}  
    & & \multicolumn{4}{p{0.6\linewidth}}{
        \cellcolor[rgb]{0.9,0.9,0.9}{
            \makecell[{{p{\linewidth}}}]{
                \texttt{\tiny{[P1$\langle$GM]}}
                \texttt{The word flock was guessed correctly. Your teammate's turn ended there.} \\
\texttt{Now provide another clue relating to some of your remaining team words and a list of the related target words. Remember to start your clue with 'CLUE: ', put a new line, and start your comma{-}separated list of target words with 'TARGETS: '. Notice: some words have been removed from the lists compared to previous requests.} \\
\\ 
\texttt{Your remaining team words are: dab.} \\
\texttt{Remaining words for your opponent are: bust.} \\
\texttt{Remaining distractor words are: steer, darn, stir, counter.} \\
\texttt{Remaining kill words are: .} \\
            }
        }
    }
    & \\ \\

    \theutterance \stepcounter{utterance}  
    & \multicolumn{4}{p{0.6\linewidth}}{
        \cellcolor[rgb]{0.9,0.9,0.9}{
            \makecell[{{p{\linewidth}}}]{
                \texttt{\tiny{[P1$\rangle$GM]}}
                \texttt{CLUE: MAKEUP} \\
\texttt{TARGETS: dab} \\
            }
        }
    }
    & & \\ \\

    \theutterance \stepcounter{utterance}  
    & & & \multicolumn{2}{p{0.3\linewidth}}{
        \cellcolor[rgb]{0.9,0.9,0.9}{
            \makecell[{{p{\linewidth}}}]{
                \texttt{\tiny{[GM$|$GM]}}
                \texttt{makeup} \\
            }
        }
    }
    & & \\ \\

    \theutterance \stepcounter{utterance}  
    & & & \multicolumn{2}{p{0.3\linewidth}}{
        \cellcolor[rgb]{0.9,0.9,0.9}{
            \makecell[{{p{\linewidth}}}]{
                \texttt{\tiny{[GM$|$GM]}}
                ['dab']
            }
        }
    }
    & & \\ \\

    \theutterance \stepcounter{utterance}  
    & & & \multicolumn{2}{p{0.3\linewidth}}{
        \cellcolor[rgb]{0.9,0.9,0.9}{
            \makecell[{{p{\linewidth}}}]{
                \texttt{\tiny{[GM$|$GM]}}
                {'word': 'dab', 'assignment': 'team'}
            }
        }
    }
    & & \\ \\

    \theutterance \stepcounter{utterance}  
    & & & \multicolumn{2}{p{0.3\linewidth}}{
        \cellcolor[rgb]{0.9,0.9,0.9}{
            \makecell[{{p{\linewidth}}}]{
                \texttt{\tiny{[GM$|$GM]}}
                \texttt{opponent has won} \\
            }
        }
    }
    & & \\ \\

    \theutterance \stepcounter{utterance}  
    & \multicolumn{4}{p{0.6\linewidth}}{
        \cellcolor[rgb]{0.95,0.95,0.95}{
            \makecell[{{p{\linewidth}}}]{
                \texttt{\tiny{[P2$\rightarrow$P1]}}
                \texttt{Post{-}episode Reflexion (not sent to GM):} \\
\texttt{New plan: My previous attempt failed because I provided a clue ("MAKEUP") that was too specific to a single word ("dab") without considering the broader context of the game state, leading to a turn end before all words were guessed. While "MAKEUP" is a valid association for "dab," the strategy of isolating words one{-}by{-}one with narrow clues is inefficient and risky if the partner doesn't immediately connect. More critically, I need to ensure that if I have multiple words remaining, I group them under a single, strong clue to maximize the chance of clearing the board in one turn. In this specific case, with only "dab" left, the clue "MAKEUP" was actually semantically sound, but the loss implies the partner might have guessed a distractor (like "darn" or "stir" if they misinterpreted "dab" as a verb related to cleaning) or the game ended due to a turn limit not explicitly shown but implied by the "LOSE" state. However, looking at the distractors (steer, darn, stir, counter), "dab" as a verb (to clean) could link to "darn" (mending) or "stir" (mixing) if the partner is confused. A better strategy is to use a clue that strongly evokes the *noun* form or a very specific context to avoid ambiguity. For "dab," a clue like "DANCE" (referring to the Dab dance) or "PAINT" (referring to dabbing paint) might be safer than "MAKEUP" if "MAKEUP" triggered a distractor. But the most likely error in a "LOSE" state with one word left is that the partner guessed a distractor. I will ensure future clues for "dab" are distinct from "darn" (mending) and "stir" (mixing). I will choose "DANCE" as it is a unique, pop{-}culture association for "dab" that has no semantic overlap with "steer," "darn," "stir," or "counter." If multiple words remain, I will group them under a unifying theme rather than splitting them.} \\
\\ 
\texttt{New plan: My previous attempt failed because the clue "MAKEUP" for "dab" likely led the partner to guess a distractor word (possibly "darn" or "stir" due to ambiguous verb associations) or the turn ended prematurely. To prevent this, I will select clues that have a unique, unambiguous association with the target word and zero overlap with the current distractor list. For "dab," instead of "MAKEUP" (which can be associated with cleaning/mending), I will use "DANCE" (referencing the specific dance move) or "PAINT" (referencing the action of applying paint), ensuring the link is strong and exclusive. If multiple words} \\
            }
        }
    }
    & & \\ \\

    \theutterance \stepcounter{utterance}  
    & & & \multicolumn{4}{p{0.6\linewidth}}{
        \cellcolor[rgb]{0.95,0.95,0.95}{
            \makecell[{{p{\linewidth}}}]{
                \texttt{\tiny{[P1$\rightarrow$P2]}}
                \texttt{Post{-}episode Reflexion (not sent to GM):} \\
\texttt{New plan: The previous attempt failed because I selected "flock," which, while related to birds, was not the intended answer for this specific instance, leading to a "LOSE" outcome despite correct formatting. The strategy of selecting a single word was correct, but the semantic choice was insufficient. Next time, I will analyze the provided list for the most direct and exclusive association with the target word 'bird'. I will prioritize words that are unique to the concept of birds (like "flock" or "dab" if it refers to a bird action) over general verbs or nouns that could apply to other contexts. If multiple candidates exist, I will choose the one with the strongest, most unambiguous link to the target, ensuring the single word selected is the definitive answer the system expects.} \\
            }
        }
    }
     \\ \\

\end{supertabular}
}

\end{document}
